\documentclass[a4paper,12pt]{article}

% --- Encoding ---
\usepackage[utf8]{inputenc}
\usepackage[T1]{fontenc}
\usepackage[english]{babel}

% --- Math ---
\usepackage{mathtools}
\usepackage{amssymb}

% --- PythonTeX ---
\usepackage{pythontex}
% Optional: enable stderr output for debugging
% \setpythontexoutputdir{.}

% --- Typography ---
\usepackage[margin=2.5cm]{geometry}
\usepackage{microtype}
\usepackage{booktabs}

% --- Graphics (for matplotlib figures) ---
\usepackage{graphicx}
\usepackage{xcolor}

% --- Navigation ---
\usepackage{hyperref}
\usepackage{cleveref}

% ============================================================
% COMPILATION:
%   pdflatex document.tex
%   pythontex document.tex      (or: python -m pythontex document.tex)
%   pdflatex document.tex
%
% Or add to .latexmkrc:
%   add_cus_dep('pytxcode', 'pytxminted', 0, 'pythontex');
%   sub pythontex { system("pythontex \"$_[0]\""); }
% ============================================================

\title{PythonTeX Document}
\author{Author Name}
\date{\today}

% ------ Define Python constants here (once, reused everywhere) ------
\begin{pycode}
# Global constants — available throughout the document
SPEED_OF_LIGHT = 299_792_458     # m/s
PLANCK_H       = 6.626e-34       # J·s
AVOGADRO       = 6.022e23        # mol⁻¹
PI             = 3.14159265358979

# Helper: format large numbers with commas
def fmt(n, decimals=2):
    return f"{n:,.{decimals}f}"
\end{pycode}

\begin{document}

\maketitle

\section{Using Python Variables}

The speed of light is $c = \py{SPEED_OF_LIGHT:,}$ m/s.

Avogadro's number is $N_A = \py{f"{AVOGADRO:.4e}"}$ mol$^{-1}$.

\section{Inline Calculations}

$\sqrt{2} = \py{2**0.5:.6f}$

$10! = \py{__import__('math').factorial(10)}$

\section{SymPy Symbolic Math}

\begin{sympycode}
from sympy import *
x = symbols('x')
f = x**4 - 3*x**2 + 2
deriv = diff(f, x)
antideriv = integrate(f, x)
\end{sympycode}

Given $f(x) = \sympy{f}$:
%
\begin{align}
  f'(x)                    &= \sympy{deriv} \\
  \int f(x)\,\mathrm{d}x   &= \sympy{antideriv} + C
\end{align}

Roots of $f(x) = 0$: $x \in \sympy{solve(f, x)}$.

\section{Generated Table}

\begin{pycode}
import math

headers = ["$n$", "$n!$", "$\ln(n!)$"]
rows = [(n, math.factorial(n), math.log(math.factorial(n))) for n in range(1, 9)]

print(r"\begin{table}[htbp]")
print(r"  \centering")
print(r"  \caption{Factorials and their logarithms.}")
print(r"  \begin{tabular}{rrl}")
print(r"    \toprule")
print(f"    {' & '.join(headers)} \\\\ \\midrule")
for n, fact, logfact in rows:
    print(f"    {n} & ${fact:,}$ & ${logfact:.4f}$ \\\\")
print(r"    \bottomrule")
print(r"  \end{tabular}")
print(r"\end{table}")
\end{pycode}

\section{Matplotlib Figure}

\begin{pycode}
import matplotlib
matplotlib.use('pgf')
import matplotlib.pyplot as plt
import numpy as np

x = np.linspace(0, 2*np.pi, 300)
fig, ax = plt.subplots(figsize=(5, 2.8))
ax.plot(x, np.sin(x), label=r'$\sin(x)$')
ax.plot(x, np.cos(x), '--', label=r'$\cos(x)$')
ax.set_xlabel(r'$x$')
ax.set_ylabel(r'$y$')
ax.legend(loc='upper right')
ax.grid(True, alpha=0.3)
fig.tight_layout()
fig.savefig('sincos.pdf')
plt.close()
\end{pycode}

\begin{figure}[htbp]
  \centering
  \includegraphics[width=0.7\linewidth]{sincos.pdf}
  \caption{Sine and cosine functions generated by Python.}
  \label{fig:sincos}
\end{figure}

\end{document}
